\section{Introduction}\label{sec:intro}

Organisations now run video through their information systems continuously. Contact centres record customer calls with a camera on, tele-health platforms stream consultations, retailers and advertisers test creative material on viewer panels, and live-commerce hosts sell to audiences in real time. Deep video models have made it possible to recognise events and activities in these streams \citep{mao2024}, and edge--cloud architectures have made it possible to do so at scale \citep{ananthanarayanan2017,gong2025}. A growing share of the value that organisations expect from such streams is affective: whether a viewer is engaged, whether a customer is frustrated, whether a presenter conveys warmth. Emotion and engagement cues derived from video have been shown to predict online engagement \citep{stappen2022}, viewer responses to advertising \citep{antonov2024}, and, through measures of empathy, firm outcomes \citep{cui2025}. In information systems (IS) research, discrete emotions expressed by consumers are known to shape how others evaluate their contributions \citep{yin2014}. Affective video analytics is therefore becoming a component of decision support, customer analytics and adaptive services \citep{picard1997,poria2017}.

Most of the affective computing literature treats recognition as a classification problem on a fixed, curated test set. Multimodal models that combine face, voice and language now recognise posed and spontaneous emotions well on such benchmarks \citep{tzirakis2017,zadeh2017,tsai2019,wagner2023}. Run on a live stream, these models raise three problems that benchmark accuracy does not capture.

The first problem is \emph{cost}. A stream must be analysed within a compute and latency budget, and the most accurate models are also the most expensive. Video analytics systems have long used cascades and adaptive configuration to cut cost \citep{kang2017,jiang2018}, but affective pipelines have rarely been designed or evaluated this way. The second problem is \emph{silent degradation}. Real streams change: a microphone picks up background noise, a bandwidth drop forces a lower bit rate, a room goes dark. Ground-truth emotion labels are not available in production, so the system cannot measure its own accuracy. The usual fallback is to monitor the model's confidence \citep{hendrycks2017}. Neural classifiers are known to become over-confident under distribution shift \citep{ovadia2019}, but how this plays out in multimodal affect pipelines, where one channel may degrade while the other stays intact, has not been examined. The third problem is \emph{accountability}. Affective inferences about people are contested on scientific grounds \citep{barrett2019}, they may differ in quality across demographic groups \citep{xu2020,mehrabi2021}, and the European Union's Artificial Intelligence Act restricts and regulates emotion recognition \citep{euaiact2024}. Organisations that deploy these systems need calibrated probabilities, group-level audits and explanations of what drove a decision \citep{rai2020,berente2021}.

These problems are design problems, so we follow the design science research paradigm \citep{hevner2004,peffers2007}. We ask: \emph{How should a multimodal affect recognition component of a streaming video information system be designed so that it remains accurate for unseen people, stays within a compute budget, detects and responds to stream degradations without labels, and supports accountability?} We break this down into four research questions:
\begin{itemize}
\item[\textbf{RQ1}] How much do audio and visual foundation-model features, and their fusion, contribute to subject-independent emotion recognition, and at what compute cost?
\item[\textbf{RQ2}] Can a confidence-gated cascade recover the accuracy of the full multimodal model at a fraction of its cost, and how should the cheap first stage be chosen?
\item[\textbf{RQ3}] How do realistic stream degradations affect each modality, and can output confidence or input-space statistics detect them without labels?
\item[\textbf{RQ4}] Can label-free monitors be turned into a gating policy that protects accuracy under degradation, how does such a policy interact with robust training, and what are the calibration and group-fairness properties of the resulting system?
\end{itemize}

We answer these questions over four build-and-evaluate cycles. We instantiate the artifact with four feature extractors of increasing cost: hand-crafted prosody, a self-supervised speech transformer (wav2vec~2.0), a mobile convolutional network, and a vision--language model (CLIP). We evaluate it on all 1440 audio-visual speech clips of the RAVDESS corpus \citep{livingstone2018} with a subject-independent protocol, measure CPU cost per clip, and inject eight degradation conditions into the audio and video streams.

The study makes four contributions. First, it reports the cost of each step of a multimodal affect pipeline alongside its accuracy. Late fusion of audio and visual foundation features raises accuracy by 10.5 percentage points over the best single modality. A cascade that runs the speech model first cuts compute by 48\% at a cost of 2.6 points of accuracy. An integrated variant that also monitors both channels saves 21\% at a cost of 1.5 points. Second, it shows a failure that matters for deployment. Under additive audio noise, clean-trained models become \emph{more} confident while they become less accurate. A cascade driven by confidence then stops escalating, and a confidence-based drift monitor stays silent. Third, it shows that cheap, per-modality monitors on the input features detect every degradation we tested, typically within 15 utterances (under a minute of speech). Each degradation is attributed to the correct channel. The same monitors also react to natural speaker changes. Fourth, it shows how monitoring, gating and robust training interact. Gating on stream-level alarms protects clean-trained models, but it harms models trained on degraded data unless it responds only to shifts outside their training envelope. These findings are condensed into six design principles for affective video information systems.

The remainder of the paper is organised as follows. Section~\ref{sec:background} reviews related work. Section~\ref{sec:approach} describes the research approach and design requirements, and Section~\ref{sec:artifact} the artifact. Section~\ref{sec:evaluation} presents the evaluation design and Section~\ref{sec:results} the results. Section~\ref{sec:discussion} discusses implications and limitations, and Section~\ref{sec:conclusion} concludes.
